\section{Exact quartics and their computational verification}
\label{sec:quartic-artifacts}

The preceding anchor formula is applied to the actual retained residual
system, not only to an informal count. The old private anchor uses two
scaling coordinates $a_i$ with $Ca_i=f_i$, and, for each $i=0,1$, the
seven-witness equations
\begin{align*}
 b_{i0}^2-2b_{i1}^2-3b_{i2}^2+6b_{i3}^2&=1,\\
 (a_i-2b_{i0})^2-8c_{i1}^2-12c_{i2}^2+24c_{i3}^2&=4.
\end{align*}
The norm-one trace-sum theorem for the quaternion algebra $(2,3)$
identifies this condition with $a_i\in R$; its ramification set is
$\{2,3\}$ and it splits at infinity
\cite[Proposition 6]{Koenigsmann2016}. The second equation has scalar
part $(a_i-2b_{i0})/2$, so the trace convention agrees exactly.
The retained constant is $C=3^{2^{24}}\ne0$.

\begin{proposition}\label{prop:computed-quartics}
There are exact integer quartics $P_{12493}$ and $P_{12465}$ with
the same rational input projection as the released source. Their
statistics are
\[
\begin{array}{c|r|r|r|r}
 &\text{total variables}&\text{residuals}&\text{monomials}&\log_2 H\ \text{bound}\\\hline
 P_{12493}&12493&11936&135667&<80\\
 P_{12465}&12465&11908&130427&<3247.
\end{array}
\]
The height of $P_{12493}$ equals the height of the release.
Both complete expanded integer coefficient tables are supplied.
\end{proposition}
\begin{proof}
For the first polynomial replace the six old anchor equations by
the three equations of Corollary~\ref{cor:five-anchor}. Their sixteen
private coordinates occur nowhere else. The five new coordinates
retain quadratic residuals and have exactly the same projection at the
interface $(C,f_0,f_1)$. Eleven variables and three residuals are removed.
The sum of the squares of all retained residuals is a quartic.

For the second polynomial first perform twenty-eight additional exact
eliminations and thirteen invertible row operations. An elimination
uses a row $cY+g=0$, with $c\in\Q^\times$ and $g$ independent of
$Y$, substitutes $Y=-g/c$ into every occurrence, and removes that row
and coordinate. Both rational solution maps are immediate. A row
operation replaces $f$ by $af+bg$ while retaining $g$, with
$a\in\Q^\times$; again both solution sets agree. The certificate
records every affected row, all denominator clearings, and the variable
remapping. Every resulting residual is still quadratic. Apply the same
eleven-variable anchor replacement afterward.

The independent checker replays the pinned baseline certificate, checks
every extra operation using exact fractions without importing the
producer, checks the literal six anchor rows and their full incidence,
verifies the nonzero defining chain of $C$, and expands every square
by ordered products. It compares every coefficient with the supplied
table and verifies that each reported variable occurs. This gives the
displayed finite statistics. The semantic anchor equivalence is proved
above; the checker verifies its exact placement. Neither part uses an
unproved arithmetic optimization of the source.
\end{proof}

The small-coefficient package provides a useful alternative when file
size matters. Under the requested variable-first ordering,
$P_{12465}$ is preferable despite its coefficient increase. Its largest
coefficient has 978 decimal digits; the exact integer is recorded in the
verification receipt, rather than reproduced here.

\subsection{Removing constants without expanding their decimal digits}

There are still fifteen fixed coordinates in $P_{12465}$, beginning
with $c_0=3^{1024}$ and continuing by $c_{k+1}=c_k^2$ for
$0\leq k<14$. Thus $c_{14}=C=3^{2^{24}}$. The new anchor coordinate
$D_0=C^2$ is fixed as well. Substitute these sixteen constants into
every remaining residual and delete their defining rows. This cannot
increase degree. The only remaining large constant is
\[
 M=C^4=3^{2^{26}},
\]
in the anchor residual
\[
 u^2+6v^2-2w^2-3r^2+2M.
\]
Here $M$ is a fixed integer, not a variable.

\begin{proposition}\label{prop:constant-quartic}
The constant substitution gives a quartic in exactly 12,449 total
variables, with 11,892 residuals and 130,378 monomials, having the
same rational input projection. Its primitive integer coefficient
height is exactly
\begin{equation}\label{eq:constant-height}
 H=4\cdot3^{134217728}+H_0,
\end{equation}
where
\begin{align*}
H_0={}&373684652\cdot10^{144}\\
&+ 355249371111055415460965317040747589616777785974\cdot10^{96}\\
&+ 770831188591770167322175457505857070439628239726\cdot10^{48}\\
&+ 443171972269170621942059432904883611120549682273.
\end{align*}
Its coefficients are completely specified by a table over $\Z[M]$
of degree at most two in the fixed constant $M$.
\end{proposition}
\begin{proof}
Each deleted coordinate is uniquely fixed by its defining equation.
Direct incidence checking shows that, outside the deleted chain, these
coordinates occur only in the displayed anchor places. Thus substitution
and reconstruction give inverse maps on the retained coordinates.
The exact expansion over $\Z[M]$ gives the stated counts.
No nonconstant monomial disappears on specializing the fixed constant:
its coefficient is $a+bM$ with $|a|<2^{3247}<M$, so it is nonzero
when $b\ne0$; a nonzero constant coefficient remains nonzero as well.
Only the coefficient of the constant monomial has an $M^2$ term.
The constant coefficient is $4M^2+H_0$. Every other coefficient has
absolute value at most $2^{3247}+24M$, which is strictly smaller.
The coefficient of $u^4$ is one, so the integer content is one and the
degree is four. Each remaining variable occurs in a residual, and a
sum of squares of real polynomials cannot lose its dependence on such
a variable: its leading coefficient in that variable is a nonzero
sum of squares. This verifies the exact variable count.
\end{proof}

The integer in \eqref{eq:constant-height} is intentionally not converted
to a decimal string. This is an exact polynomial specification by an
ordinary integer power, not a numerical approximation. An elementary
upper bound is $H<2^{268435459}$. The smaller variable count is purchased
with very large coefficients, in accordance with the specified priority.

\subsection{Verification boundaries}

The files separate producer output from independent verification. The
constant table is checked as a polynomial over $\Z[M]$, so there is no
floating-point arithmetic or gigantic decimal conversion in that check.
For the preceding elimination suffix, seven deliberate mutations of
pivots, solved numerators, affected-row lists, substitution patches,
and row operations were rejected. These checks test the advertised
finite transformations.
They neither test all rational inputs nor certify the manuscript's
restricted converse, height estimate, or prime-pattern argument.
